%% kapa: exam=1 round=26 year=2015 period=2 form=A subject=accounting subject_ko=회계학 seq=39 number=79
%% source: Q-Net 감정평가사 기출문제 (cst00309) artlSeq=5138001 file=extracted/1008/5138001/제26회_감정평가사 1차_기출/제26회_감정평가사 1차_기출/2교시_감평_A형.pdf
%% conversion: kordoc 4.24.51 -> tools/gen1.py (auto); review=proofread-sonnet; flags=none
\begin{question}{79}
표준원가계산제도를 채택하고 있는 (주)감평의 20×5년 6월의 제조간접원가 예산과 실제발생액은 다음과 같다.

{\small\centering\begin{tabular}{|l|r|r|}\hline
 & \multicolumn{1}{c|}{예산} & \multicolumn{1}{c|}{실제발생액} \\ \hline
변동제조간접원가: & & \\
\hspace{1.5em}간접재료원가 & ￦4,480 & ￦4,310 \\
\hspace{1.5em}간접노무원가 & 3,360 & 3,630 \\
고정제조간접원가: & & \\
\hspace{1.5em}공장임대료 & 10,640 & 10,630 \\
\hspace{1.5em}감가상각비 & 14,280 & 14,260 \\ \hline
총 제조간접원가 & ￦32,760 & ￦32,830 \\ \hline
\end{tabular}\par}
\medskip
(주)감평은 고정제조간접원가의 예산을 월간 2,800기계시간을 기준으로 수립하였다. 20×5년 6월 실제 기계시간은 2,730시간이다. 만약 실제 생산량에 허용된 표준 기계시간이 2,860시간이라면, 6월의 고정제조간접원가 조업도차이는?
\begin{choices}
  \item ￦534 불리
  \item ￦623 불리
  \item ￦702 불리
  \item ￦534 유리
  \item ￦623 유리
\end{choices}
\end{question}
\answer{4}
